Fix \(p,m,c\), put \(h=m-c\) and \(q=m-2c=h-c\), and fix rational thresholds \(\gamma _0,\zeta _0,\lambda _0>0\) and \(\overline\kappa,\overline M\geq1\).  Define
\[
 L_0=\bigl(p!\,\overline\kappa^{\,p-1}\bigr)^{1/p},\qquad
 K_0=\frac{6\overline M}{\gamma _0}
 \sqrt{p(p-1)(1+L_0^2)},
\]
\[
 \rho _0=\min\left\{1,\frac{1}{2K_0\overline M}\right\},
 \qquad C_0=16K_0L_0^3.
\]
These are effectively represented positive real algebraic numbers.

Let \(\mathfrak T\) be the finite union, over \(H_1,H_2\subseteq E\) with \(|H_1|=|H_2|=h\) and \(S\subseteq H_1\cap H_2\) with \(|S|=q\), of tuples
\[
 \tau=\bigl(H_1,H_2,S,D,U,\Omega,(s_e,\Sigma_e)_{e\in H_1},
 A,V,\Psi,(t_e,\Gamma_e)_{e\in H_2}\bigr)
\]
where \(D,U,A,V\in\mathbb R^{p\times p}\), the matrices \(\Omega,\Psi,\Sigma_e,\Gamma_e\) are symmetric \(p\times p\) matrices, and \(s_e,t_e\in\mathbb R^p\), satisfying the following closed constraints.  The inverse variables obey \(DU=UD=I_p\) and \(AV=VA=I_p\).  Both \(D\) and \(A\) have unit diagonal; for every directed simple cycle \(C\), the corresponding products in \(I_p-D\) and \(I_p-A\) have absolute value at most \(1-\zeta _0\); and
\[
 \|D\|_{\mathrm{op}}\|U\|_{\mathrm{op}}\leq\overline\kappa,
 \qquad
 \|A\|_{\mathrm{op}}\|V\|_{\mathrm{op}}\leq\overline\kappa.
\]
The nuisance matrices, covariance matrices, and shifts satisfy
\[
 \Omega,\Psi\succeq0,qquad
 \Sigma_e,\Gamma_e\succeq\lambda _0I_p,qquad
 s_e,t_e\in[0,\infty)^p,
\]
\[
 D\Sigma_eD^{\mathsf T}=\Omega+\operatorname{diag}(s_e)\quad(e\in H_1),
 \qquad
 A\Gamma_eA^{\mathsf T}=\Psi+\operatorname{diag}(t_e)\quad(e\in H_2).
\]
Thus \(\lambda_0\) audits covariance interiority only and does not fix the additive diagonal gauge of either invariant nuisance. Each of the two shift designs has affine separation at least \(\gamma _0\), with the separation computed using all subsets of its honest set having cardinality \(q\), and each side satisfies its scale bound
\[
 1+\|\Omega\|_{\mathrm{op}}+
 \max_{e\in H_1}\|\Sigma_e\|_{\mathrm{op}}+
 \max_{e\in H_1}\|\operatorname{diag}(s_e)\|_{\mathrm{op}}\leq\overline M,
\]
\[
 1+\|\Psi\|_{\mathrm{op}}+
 \max_{e\in H_2}\|\Gamma_e\|_{\mathrm{op}}+
 \max_{e\in H_2}\|\operatorname{diag}(t_e)\|_{\mathrm{op}}\leq\overline M.
\]
All operator-norm, eigenvalue, and finite minimum--maximum conditions here are interpreted through their equivalent first-order real-polynomial formulas.  Define the full covariance-equation residual on the overlap by
\[
 \mathfrak r(\tau)=
 \max_{e\in S}
 \left\|A\Sigma_eA^{\mathsf T}-\Psi-\operatorname{diag}(t_e)\right\|_{\mathrm F},
\]
and let \(\mathfrak T_{\mathrm{far}}\) be the subfamily satisfying
\[
 \|AU-I_p\|_{\mathrm F}\geq\rho _0.
\]
Set \(\mu _0=1\) if \(\mathfrak T_{\mathrm{far}}=\varnothing\), and otherwise set
\[
 \mu _0=\min_{\tau\in\mathfrak T_{\mathrm{far}}}\mathfrak r(\tau).
\]
Starting with \(j=1,2,\ldots\), use effective real quantifier elimination to decide whether some \(\tau\in\mathfrak T_{\mathrm{far}}\) has \(\mathfrak r(\tau)<2^{-j}\), and let \(\underline\mu=2^{-j}\) at the first negative answer.  Finally define
\[
 r_0=\frac{\underline\mu}{2\sqrt p\,L_0^2}.
\]
The package consisting of \(L_0,K_0,\rho _0,C_0,\mathfrak T,\mu _0,\underline\mu,r_0\) is \(\operatorname{AuditedReleaseContraction}(p,m,c,\gamma _0,\zeta _0,\lambda _0,\overline\kappa,\overline M)\).  Positivity and termination of the search are theorem properties, not part of this definition.